% Краткое сообщение для журнала «Доклады Российской академии наук.
% Математика, информатика, процессы управления».
% Оформление --- официальный стилевой файл DAN2E.STY (DanMath_for_Authors.zip,
% sciencejournals.ru/journal/danmiup/), здесь --- копия dan2e.sty.
% Правила (guid.pdf, сверено 14.09.2026): <= 25 000 знаков с таблицами и
% библиографией (<= 25 источников), представление академиком РАН, аннотация
% <= 1000 знаков без ссылок, 3--10 ключевых слов, метаданные на русском и
% английском, ссылки на неопубликованные работы не допускаются,
% самоцитирование <= 30 %, одновременная подача в другой журнал запрещена.
%
% Содержание намеренно НЕ пересекается с заметкой УМН (note-umn/note-2p.tex):
% там --- четыре явные конструкции с точными сертификатами (43/132/1029/7203),
% здесь --- аналитическая часть: тождество для эйзенштейновых решёток,
% продуктовое правило, оценки 45619 (R^10), 4*7^12 (R^25), 19*7^12 (R^26)
% и оболочечный пол индекса (2401 неулучшаемо на E_8). Все доказательства
% полные и не опираются на вычисления и на неопубликованные работы.
\documentclass{article}

\usepackage[T2A]{fontenc}
\usepackage[utf8]{inputenc}
\usepackage[russian]{babel}
\usepackage{amsmath}
\usepackage{amssymb}
\usepackage{amsthm}
\usepackage{mathrsfs}

\usepackage{dan2e}

\theoremstyle{plain}
\newtheorem{theorem}{Теорема}
\newtheorem{proposition}{Предложение}
\newtheorem{lemma}{Лемма}
\newtheorem{corollary}{Следствие}
\theoremstyle{remark}
\newtheorem{remark}{Замечание}

\newcommand{\R}{\mathbb{R}}
\newcommand{\Z}{\mathbb{Z}}
\newcommand{\CC}{\mathbb{C}}
\newcommand{\diam}{\operatorname{diam}}
\newcommand{\dist}{\operatorname{dist}}
\newcommand{\vol}{\operatorname{vol}}

\begin{document}

% Том, год и страницы проставляет редакция.
\Volume{000}
\Year{2026}
\Pages{000--000}

% TODO: сверить УДК с редакцией.
\udk{514.174.5+519.174}

\title{Точные ширины решёточных раскрасок евклидовых пространств:
тождество для эйзенштейновых решёток, продуктовое правило и оценка
$\chi(\mathbb R^{10})\le45619$}

% Порядок авторов --- по вкладу (решение от 22.08.2026).
\author{Л.\,Л.~Иванов\Addressmark[1]\Emailmark[1],
Н.\,В.~Глушкова\Addressmark[2]\Emailmark[2]}

% TODO: место работы обоих авторов (полное название учреждения, город).
\Addresstext[1]{Москва, Россия}
\Addresstext[2]{Московский физико-технический институт (национальный
исследовательский университет), Долгопрудный, Московская обл., Россия}

\Emailtext[1]{leo.ivanov@gmail.com}
\Emailtext[2]{glushkova.nv@phystech.edu}

\markboth{Л.\,Л.~Иванов, Н.\,В.~Глушкова}{Точные ширины решёточных
раскрасок евклидовых пространств}

% TODO: представление действительного члена РАН (п. 1 правил).
\presentedby{Представлено академиком РАН \ldots}

\dateA{\ldots}

\alttitle{Exact widths of lattice colorings of Euclidean spaces:
an identity for Eisenstein lattices, the product rule, and the bound
$\chi(\mathbb R^{10})\le45619$}

\altauthor{L.\,L.~Ivanov\Addressmark[a]\Emailmark[1],
N.\,V.~Glushkova\Addressmark[b]\Emailmark[2]}

\altAddresstext[a]{Moscow, Russian Federation}
\altAddresstext[b]{Moscow Institute of Physics and Technology (National
Research University), Dolgoprudny, Moscow Region, Russian Federation}

\altpresentedby{Presented by Academician of the RAS \ldots}

\maketitle

\begin{abstract}
Рассматриваются раскраски пространства $\R^n$, при которых никакие две
одноцветные точки не находятся на расстоянии из отрезка $[1,\ell]$; при
$\ell=1$ наименьшее число цветов есть классическое хроматическое число
$\chi(\R^n)$. Для периодической раскраски ячеек Вороного решётки $\Lambda$
классами подрешётки $\Gamma$ ширина запрещённого интервала равна отношению
расстояния между одноцветными ячейками к диаметру ячейки. Доказано
тождество: для любой решётки $\Lambda$ с эйзенштейновой структурой расстояние
между одноцветными ячейками подрешётки $(3+\omega)\Lambda$ равно
$\sqrt{7/3}\,\lambda_1(\Lambda)$, где $\lambda_1$~--- длина кратчайшего
вектора; отсюда точные ширины всех известных раскрасок в $7^{n/2}$ цветов.
Доказано продуктовое правило: ширина ортогонального произведения раскрасок
с ширинами $d_i$ при оптимальном масштабировании блоков равна
$(\sum_i d_i^{-2})^{-1/2}$. Как следствие, $\chi(\R^{10})\le45619$~---
первая оценка в $\R^{10}$ ниже классической $3^{10}$, а также
$\chi(\R^{25})\le4\cdot7^{12}$ и $\chi(\R^{26})\le19\cdot7^{12}$. Доказано,
что среди подрешёток решётки $E_8$, задающих правильную раскраску, индекс
$2401$ неулучшаем.
\end{abstract}

\begin{keywords}
хроматическое число, проблема Нельсона--Хадвигера, запрещённый интервал
расстояний, решётка, ячейка Вороного, радиус покрытия, эйзенштейновы
решётки, решётка $E_8$, постоянная Эрмита
\end{keywords}

\begin{altabstract}
We consider colorings of $\R^n$ in which no two points of the same color
are at a distance from the segment $[1,\ell]$; for $\ell=1$ the minimum
number of colors is the classical chromatic number $\chi(\R^n)$. For a
periodic coloring of the Voronoi cells of a lattice $\Lambda$ by the cosets
of a sublattice $\Gamma$, the width of the forbidden interval is the ratio
of the distance between same-colored cells to the diameter of a cell. We
prove an identity: for every lattice $\Lambda$ with an Eisenstein structure,
the distance between same-colored cells of the sublattice $(3+\omega)\Lambda$
equals $\sqrt{7/3}\,\lambda_1(\Lambda)$, where $\lambda_1$ is the length of a
shortest vector; this gives the exact widths of all known colorings with
$7^{n/2}$ colors. We prove a product rule: the width of an orthogonal
product of colorings with widths $d_i$, under the optimal scaling of the
blocks, equals $(\sum_i d_i^{-2})^{-1/2}$. As a consequence,
$\chi(\R^{10})\le45619$, the first bound in $\R^{10}$ below the classical
$3^{10}$, and also $\chi(\R^{25})\le4\cdot7^{12}$ and
$\chi(\R^{26})\le19\cdot7^{12}$. We prove that among the sublattices of
$E_8$ defining a proper coloring the index $2401$ cannot be improved.
\end{altabstract}

\begin{altkeywords}
chromatic number, Nelson--Hadwiger problem, forbidden distance interval,
lattice, Voronoi cell, covering radius, Eisenstein lattices, the lattice
$E_8$, Hermite constant
\end{altkeywords}

\section*{1. Введение и результаты}

Через $\chi(\R^n,[1,\ell])$, $\ell\ge1$, обозначим наименьшее число цветов
раскраски пространства $\R^n$, при которой никакие две одноцветные точки не
находятся на расстоянии из отрезка $[1,\ell]$. При $\ell=1$ это классическое
хроматическое число $\chi(\R^n)$ (проблема Нельсона--Хадвигера, см.~\cite{Soifer,deGrey}); интервальная постановка изучалась в~\cite{Ivanov2011}.
Очевидно, $\chi(\R^n)\le\chi(\R^n,[1,\ell])$.

Все известные верхние оценки $\chi(\R^n)$ в малых размерностях получены
\emph{решёточными раскрасками}. Пусть $\Lambda\subset\R^n$~--- решётка,
$\Gamma\subset\Lambda$~--- подрешётка индекса $k=[\Lambda:\Gamma]$,
$V_0$~--- ячейка Вороного нуля (выпуклый центрально-симметричный
многогранник, $\diam V_0=2R$, где $R$~--- радиус покрытия $\Lambda$);
ячейку $a+V_0$, $a\in\Lambda$, окрасим цветом класса $a+\Gamma$. Положим
\begin{equation}\label{eq:defs}
D(v)=\dist\bigl(V_0,\,v+V_0\bigr),\quad
D(\Gamma)=\min_{v\in\Gamma\setminus\{0\}}D(v),\quad
d(\Lambda,\Gamma)=\frac{D(\Gamma)}{\diam V_0},\quad
\rho(\Lambda)=\frac{\diam V_0}{\lambda_1(\Lambda)},
\end{equation}
где $\lambda_1(\Lambda)$~--- длина кратчайшего ненулевого вектора. Две
одноцветные точки лежат либо в одной ячейке (расстояние $\le\diam V_0$),
либо в двух ячейках, центры которых различаются на $v\in\Gamma\setminus\{0\}$
(расстояние $\ge D(\Gamma)$). Выбрав масштаб $s$ с $s\diam V_0<1<\ell<sD(\Gamma)$,
что возможно в точности при $\ell<d(\Lambda,\Gamma)$, получаем
\begin{equation}\label{eq:crit}
\chi(\R^n,[1,\ell])\le[\Lambda:\Gamma]\qquad\text{при всех }\ell<d(\Lambda,\Gamma).
\end{equation}
Величину $d(\Lambda,\Gamma)$ называем \emph{шириной} раскраски; раскраска
\emph{правильна}, если $d\ge1$. Лучшие известные оценки в размерностях
$4$--$10$ и $24$ принадлежат работе~\cite{ABPR}; в размерностях
$n=4,6,8,24$ они получены на решётках $D_4$, $E_6^*$, $E_8$ и решётке Лича
$\Lambda_{24}$, снабжённых структурой модуля над кольцом Эйзенштейна
$\Z[\omega]$, $\omega=e^{2\pi i/3}$, с подрешёткой $\Gamma=(3+\omega)\Lambda$
индекса $|3+\omega|^n=7^{n/2}$; в~\cite{ABPR} доказана лишь правильность
этих раскрасок. Ниже решётка называется \emph{эйзенштейновой}, если на ней
задано умножение на $\omega$, являющееся изометрией; тогда $\Lambda$~---
$\Z[\omega]$-модуль, и для $\alpha\in\Z[\omega]\setminus\{0\}$ подрешётка
$\alpha\Lambda$ имеет индекс $|\alpha|^n$. Через $H_1\subset\CC$ обозначим
ячейку Вороного решётки $\Z[\omega]\subset\CC=\R^2$: правильный
шестиугольник с фасетами на расстоянии $1/2$ от нуля и вершинами на
расстоянии $1/\sqrt3$.

\begin{theorem}\label{thm:eis}
Пусть $\Lambda\subset\R^n$~--- эйзенштейнова решётка. Тогда для любого
$\alpha\in\Z[\omega]\setminus\{0\}$
\begin{equation}\label{eq:alpha}
D(\alpha\Lambda)=2\,\lambda_1(\Lambda)\cdot\dist\bigl(\alpha/2,\,H_1\bigr),
\end{equation}
причём при $|\alpha|>1$ минимум в~\eqref{eq:defs} достигается ровно на
векторах $\alpha w$ с кратчайшими $w$. В частности,
\begin{equation}\label{eq:eis}
D\bigl((3+\omega)\Lambda\bigr)^2=\tfrac73\,\lambda_1(\Lambda)^2,\qquad
d\bigl(\Lambda,(3+\omega)\Lambda\bigr)=\frac{\sqrt{7/3}}{\rho(\Lambda)},
\end{equation}
а для целого $m\ge2$ выполнено $D(m\Lambda)=(m-1)\lambda_1(\Lambda)$.
\end{theorem}

\begin{corollary}\label{cor:widths}
Для $D_4$, $E_6^*$, $E_8$ и $\Lambda_{24}$ имеем $\rho=\sqrt2$, поэтому
ширина каждой из четырёх раскрасок~\cite{ABPR} в $7^{n/2}$ цветов равна
ровно $\sqrt{7/6}$:
$\chi(\R^n,[1,\ell])\le7^{n/2}$ при $\ell<\sqrt{7/6}$ и $n=4,6,8,24$, и в
пределах этих конструкций $\sqrt{7/6}$ не улучшить. Для решётки
Коксетера--Тодда $K_{12}$ имеем $\rho=\sqrt{8/3}$, так что
$d(K_{12},(3+\omega)K_{12})=\sqrt{7/8}<1$: раскраска в $7^6$ цветов в $\R^{12}$
этим путём не получается; зато $d(K_{12},3K_{12})=\sqrt{3/2}$, то есть
$\chi(\R^{12},[1,\ell])\le3^{12}$ при $\ell<\sqrt{3/2}$.
\end{corollary}

Пороговая константа $\sqrt{7/3}$ в~\cite[замечание~9]{ABPR} упоминается как
достаточное условие правильности ($\rho\le\sqrt{7/3}$), однако «решётки,
использующие это усиление», там названы неизвестными; тождество~\eqref{eq:eis}
показывает, что условие является и необходимым, и даёт точную ширину для
каждой решётки. Отрицательный вывод о $K_{12}$ (частичный ответ на открытый
вопрос~2 работы~\cite{ABPR}) опирается именно на верхнюю половину
тождества.

\begin{theorem}[продуктовое правило]\label{thm:prod}
Пусть $\Lambda=\bigoplus_{i=1}^{s}\Lambda_i$~--- ортогональная прямая сумма
решёток $\Lambda_i\subset\R^{n_i}$, $\Gamma=\bigoplus_i\Gamma_i$,
$\Gamma_i\subset\Lambda_i$, $d_i=d(\Lambda_i,\Gamma_i)$,
$k_i=[\Lambda_i:\Gamma_i]$. Тогда максимум ширины $d(\Lambda,\Gamma)$ по
всем независимым масштабированиям блоков равен
$\bigl(\sum_i d_i^{-2}\bigr)^{-1/2}$ и достигается при
$\diam^2V_0(\Lambda_i)\propto d_i^{-2}$. Следовательно,
$\chi(\R^n,[1,\ell])\le\prod_i k_i$, $n=\sum_i n_i$, при всех
$\ell<\bigl(\sum_i d_i^{-2}\bigr)^{-1/2}$.
\end{theorem}

\begin{theorem}\label{thm:bounds}
$\chi(\R^{10},[1,\ell])\le2401\cdot19=45619$ и
$\chi(\R^{26},[1,\ell])\le19\cdot7^{12}$ при всех $\ell<\sqrt{217/214}$;
$\chi(\R^{9},[1,\ell])\le2401\cdot4=9604$ и
$\chi(\R^{25},[1,\ell])\le4\cdot7^{12}$ при всех $\ell<\sqrt{63/61}$.
В частности, $\chi(\R^{10})\le45619$, $\chi(\R^{25})\le4\cdot7^{12}$,
$\chi(\R^{26})\le19\cdot7^{12}$.
\end{theorem}

Прежние оценки в этих размерностях~--- $3^{10}=59049$, $3^{25}$, $3^{26}$
и $17253$~\cite{ABPR}; $45619$~--- первая оценка в $\R^{10}$ ниже
классической величины $3^n$. Оценка $9604$ в $\R^9$ приведена как
кратчайший, чисто аналитический путь ниже $17253$; другой, вычислительной
конструкцией (ламинированием той же раскраски $E_8/2401$ с проверкой в
точной рациональной арифметике) авторами получена более сильная оценка
$\chi(\R^9)\le7203$, которая публикуется отдельно.

\begin{theorem}\label{thm:e8}
Пусть $\Lambda=E_8$ и подрешётка $\Gamma\subset\Lambda$ задаёт правильную
раскраску. Тогда $[\Lambda:\Gamma]\ge2401$; равенство достигается на
$\Gamma=(3+\omega)E_8$. Для объёмноцентрированной кубической решётки
$\Lambda=A_3^*$ аналогично $[\Lambda:\Gamma]\ge15$.
\end{theorem}

Общая нижняя граница Кулсона $2^{n+1}-1$~\cite{Coulson,ABPR} даёт для $E_8$
лишь $511$; теорема~\ref{thm:e8} отвечает для этой решётки на открытый
вопрос~8 работы~\cite{ABPR} о нижних границах индекса, зависящих от
дополнительных характеристик решётки.

\section*{2. Доказательство теоремы~\ref{thm:eis}}

\begin{lemma}\label{lem:mid}
$D(v)=2\dist(v/2,\,V_0)$ для всех $v\in\R^n$.
\end{lemma}

\begin{proof}
Если $x\in V_0$, $y\in v+V_0$, то $x'=v-y\in V_0$, и по выпуклости
$(x+x')/2\in V_0$, откуда $|x-y|=2\,|(x+x')/2-v/2|\ge2\dist(v/2,V_0)$.
Обратно, для ближайшей к $v/2$ точки $x\in V_0$ положим $y=v-x\in v+V_0$;
тогда $|x-y|=2|x-v/2|=2\dist(v/2,V_0)$.
\end{proof}

\begin{proposition}[планарная нижняя граница]\label{prop:planar}
Пусть $\Lambda$ эйзенштейнова, $w\in\Lambda\setminus\{0\}$,
$P=\operatorname{span}(w,\omega w)$ и $H=\|w\|H_1$~--- ячейка Вороного
плоской решётки $\Z[\omega]w\subset P$. Тогда ортогональная проекция
$\pi_P(V_0)$ содержится в $H$, и для любого $\alpha\in\Z[\omega]$
\[
\dist\bigl(\alpha w/2,\,V_0\bigr)\ \ge\ \|w\|\cdot\dist\bigl(\alpha/2,\,H_1\bigr).
\]
\end{proposition}

\begin{proof}
Векторы $\pm w$, $\pm\omega w$, $\pm(1+\omega)w=\mp\omega^2w$ лежат в
$\Lambda$ и имеют длину $\|w\|$, поскольку умножение на обратимые элементы
$\Z[\omega]$~--- изометрия. Шесть полупространств
$\{x:\langle x,u\rangle\le\|u\|^2/2\}$ по этим $u$ содержат $V_0$, а их
нормали лежат в $P$; поэтому $\pi_P(V_0)$ содержится в шестиугольнике,
высекаемом в $P$ теми же неравенствами,~--- в ячейке Вороного решётки
$\Z[\omega]w$, изометричной $\|w\|H_1$. Точка $\alpha w/2$ лежит в $P$,
проекция не увеличивает расстояний, значит
$\dist(\alpha w/2,V_0)\ge\dist_P(\alpha w/2,H)=\|w\|\dist(\alpha/2,H_1)$.
\end{proof}

\begin{lemma}[точка в ячейке]\label{lem:cell}
Пусть $u_1,\dots,u_s$~--- кратчайшие векторы решётки $\Lambda$,
$q=\sum_ic_iu_i$, $c_i\ge0$, $\sum_ic_i\le1$, и
$\langle q,u_j\rangle\le\frac12\lambda_1^2$ при всех $j$. Тогда $q\in V_0$.
\end{lemma}

\begin{proof}
Условие $x\in V_0$ равносильно $\langle x,r\rangle\le\frac12\|r\|^2$ для
всех $r\in\Lambda$. Пусть $r\in\Lambda\setminus\{0\}$. Если $r\ne u_i$, то
$\|r-u_i\|^2\ge\lambda_1^2=\|u_i\|^2$, откуда
$\langle r,u_i\rangle\le\frac12\|r\|^2$. Если $r\notin\{u_1,\dots,u_s\}$, то
$\langle q,r\rangle\le\bigl(\sum_ic_i\bigr)\frac12\|r\|^2\le\frac12\|r\|^2$;
если же $r=u_j$, требуемое неравенство есть условие леммы.
\end{proof}

\begin{proposition}[верхняя граница]\label{prop:up}
Пусть $\Lambda$ эйзенштейнова и $w$~--- её кратчайший вектор. Тогда все
шесть вершин шестиугольника $H=\|w\|H_1\subset P$ лежат в $V_0$; в
частности, $V_0\cap P=H$ и
$\dist(\alpha w/2,V_0)=\|w\|\dist(\alpha/2,H_1)$ для всех
$\alpha\in\Z[\omega]$.
\end{proposition}

\begin{proof}
Из $1+\omega+\omega^2=0$ следует $w+\omega w=-\omega^2w$, поэтому
$\|w+\omega w\|=\|w\|$, и, раскрывая квадрат,
\begin{equation}\label{eq:ip}
\langle w,\omega w\rangle=-\tfrac12\|w\|^2=-\tfrac12\lambda_1^2 .
\end{equation}
Вершины $H$~--- это $\pm\omega^kq$, $k=0,1,2$, где
$q=\frac13(2w+\omega w)$~--- центроид треугольника $\{0,w,-\omega^2w\}$;
достаточно проверить $q\in V_0$, остальные вершины получаются заменой $w$ на
$\pm\omega^kw$. Векторы $w$ и $\omega w$ кратчайшие,
$q=\frac23w+\frac13\omega w$, и по~\eqref{eq:ip}
$\langle q,w\rangle=\frac13(2-\frac12)\lambda_1^2=\frac12\lambda_1^2$,
$\langle q,\omega w\rangle=\frac13(-1+1)\lambda_1^2=0$; лемма~\ref{lem:cell}
даёт $q\in V_0$. По выпуклости $H\subseteq V_0\cap P$, а
$V_0\cap P\subseteq\pi_P(V_0)\subseteq H$ по предложению~\ref{prop:planar};
значит, $V_0\cap P=H$. Для $\alpha w/2\in P$ отсюда
$\dist(\alpha w/2,V_0)\le\dist_P(\alpha w/2,H)=\|w\|\dist(\alpha/2,H_1)$,
а обратное неравенство~--- предложение~\ref{prop:planar}.
\end{proof}

\begin{proof}[Доказательство теоремы~\ref{thm:eis}]
Всякий ненулевой вектор $\alpha\Lambda$ имеет вид $\alpha w$,
$w\in\Lambda\setminus\{0\}$, и по лемме~\ref{lem:mid} и
предложению~\ref{prop:planar}
$D(\alpha w)\ge2\|w\|\dist(\alpha/2,H_1)\ge2\lambda_1\dist(\alpha/2,H_1)$,
причём для кратчайшего $w$ по предложению~\ref{prop:up} здесь равенство.
Это даёт~\eqref{eq:alpha}. Если $|\alpha|>1$, то $|\alpha|^2\ge3$ и
$|\alpha/2|\ge\sqrt3/2>1/\sqrt3$, так что $\alpha/2\notin H_1$ и
$\dist(\alpha/2,H_1)>0$; тогда для некратчайшего $w$ первое неравенство
в цепочке строгое, и минимум достигается только на кратчайших векторах.
Далее, при $\|w\|=1$ точка $(3+\omega)/2=(5/4,\sqrt3/4)$ лежит в нормальном
конусе вершины $(1/2,\,1/(2\sqrt3))$ шестиугольника $H_1$, и квадрат
расстояния до неё равен $9/16+1/48=7/12$; отсюда~\eqref{eq:eis}. Наконец,
для вещественного $\alpha=m$ ближайшая к $m/2$ точка $H_1$~--- середина
фасеты $1/2$, и $\dist(m/2,H_1)=(m-1)/2$.
\end{proof}

\begin{proof}[Доказательство следствия~\ref{cor:widths}]
Решётки $D_4$, $E_6^*$, $E_8$, $\Lambda_{24}$ и $K_{12}$ эйзенштейновы
\cite{ABPR,ConwaySloane}. Их отношения $\rho=2R/\lambda_1$ известны точно
\cite{ConwaySloane}: $\sqrt2$ для первых четырёх ($E_8$: $\lambda_1^2=2$,
$R=1$; $\Lambda_{24}$: $\lambda_1^2=4$, $R=\sqrt2$) и $\sqrt{8/3}$ для
$K_{12}$ ($\lambda_1^2=4$, $R^2=8/3$). Остаётся подставить в~\eqref{eq:eis},
а для $3K_{12}$~--- в $d=(m-1)/\rho$ при $m=3$.
\end{proof}

\section*{3. Продуктовое правило и размерности $9$, $10$, $25$, $26$}

\begin{proof}[Доказательство теоремы~\ref{thm:prod}]
Ячейка Вороного прямой суммы есть произведение ячеек, поэтому квадрат
диаметра $V_0(\Lambda)$ равен сумме квадратов диаметров
$s_i=\diam^2V_0(\Lambda_i)$ ячеек блоков, а для
$v=(v_1,\dots,v_s)$ выполнено
\[
\dist^2(v/2,V_0)=\sum_i\dist^2\bigl(v_i/2,\,V_0(\Lambda_i)\bigr),
\qquad\text{то есть}\qquad D(v)^2=\sum_iD_i(v_i)^2
\]
по лемме~\ref{lem:mid}. Так как $D_i(0)=0$, а
$D_i(v_i)\ge D(\Gamma_i)$ при $v_i\ne0$, минимум по $\Gamma\setminus\{0\}$
достигается на векторе с единственной ненулевой компонентой:
$D(\Gamma)^2=\min_iD(\Gamma_i)^2=\min_id_i^2s_i$, и
$d(\Lambda,\Gamma)^2=\min_i\bigl(d_i^2s_i\bigr)\big/\sum_js_j$.
Масштабирование блока не меняет $d_i$, так что $s_i>0$ произвольны;
при $s_i=S/d_i^2$ все $d_i^2s_i$ равны $S$ и
$d(\Lambda,\Gamma)^2=\bigl(\sum_jd_j^{-2}\bigr)^{-1}$. При любом выборе
$s_i$ положим $t_i=d_i^2s_i$; тогда
$\sum_js_j=\sum_jt_jd_j^{-2}\ge\min_it_i\cdot\sum_jd_j^{-2}$, то есть
$d(\Lambda,\Gamma)^2\le\bigl(\sum_jd_j^{-2}\bigr)^{-1}$. Индекс $\Gamma$ в
$\Lambda$ равен $\prod_ik_i$, и остаётся~\eqref{eq:crit}.
\end{proof}

Величина $d_i^{-2}$~--- \emph{цена} блока при бюджете $1$: ширина
оказывается расходуемым ресурсом, а блоки могут иметь разный индекс.

\begin{proof}[Доказательство теоремы~\ref{thm:bounds}]
\emph{Блок $E_8$.} $\Gamma_1=(3+\omega)E_8$, индекс $7^4=2401$, по
следствию~\ref{cor:widths} $d_1^2=7/6$, цена $6/7$.

\emph{Плоский блок.} $\Lambda_2=\Z[\omega]\subset\CC$ (шестиугольная
решётка, $\lambda_1=1$, $V_0=H_1$, $\diam^2V_0=4/3$),
$\Gamma_2=(5+2\omega)\Z[\omega]$, индекс $N(5+2\omega)=25-10+4=19$. По
теореме~\ref{thm:eis} $D(\Gamma_2)=2\dist\bigl((5+2\omega)/2,H_1\bigr)$;
точка $(5+2\omega)/2=(2,\sqrt3/2)$ лежит в нормальном конусе вершины
$(1/2,\,1/(2\sqrt3))$, квадрат расстояния равен $9/4+1/3=31/12$, так что
$D(\Gamma_2)^2=31/3$, $d_2^2=(31/3)/(4/3)=31/4$, цена $4/31$.

\emph{Одномерный блок.} $\Lambda_3=\Z$, $\Gamma_3=4\Z$: ячейка~--- отрезок
длины $1$, ближайшие одноцветные ячейки $[-1/2,1/2]$ и $[7/2,9/2]$ удалены
на $3$, так что $d_3=3$, цена $1/9$.

\emph{Блок Лича.} $\Gamma_4=(3+\omega)\Lambda_{24}$, индекс $7^{12}$,
$d_4^2=7/6$, цена $6/7$.

Суммы цен: $6/7+4/31=214/217<1$ для пар $(E_8,\Lambda_2)$ и
$(\Lambda_{24},\Lambda_2)$~--- размерности $10$ и $26$, индексы $2401\cdot19$
и $7^{12}\cdot19$, ширина $\sqrt{217/214}$; $6/7+1/9=61/63<1$ для пар
$(E_8,\Z)$ и $(\Lambda_{24},\Z)$~--- размерности $9$ и $25$, индексы
$2401\cdot4$ и $7^{12}\cdot4$, ширина $\sqrt{63/61}$. Остаётся применить
теорему~\ref{thm:prod}.
\end{proof}

\begin{remark}
Блок $E_8/2401$ расходует $6/7$ бюджета, поэтому в произведении может быть
лишь один такой блок, а второй блок обязан иметь $d\ge\sqrt7$.
Одномерный блок из трёх цветов ($\Gamma=3\Z$, $d=2$) не проходит:
$6/7+1/4>1$. Численный перебор плоских решёточных раскрасок показывает, что
ширина $\sqrt7$ впервые достигается при индексе $19$.
\end{remark}

\section*{4. Пол индекса на фиксированной решётке}

\begin{lemma}[инрадиусная граница]\label{lem:inr}
$D(v)\le\|v\|-\lambda_1(\Lambda)$ при $\|v\|\ge\lambda_1(\Lambda)$.
\end{lemma}

\begin{proof}
Шар радиуса $\lambda_1/2$ с центром в нуле содержится в $V_0$, поэтому
$\dist(v/2,V_0)\le\|v\|/2-\lambda_1/2$; остаётся лемма~\ref{lem:mid}.
\end{proof}

\begin{proposition}[оболочечный пол]\label{prop:shell}
Пусть $s^*>0$ таково, что $D(v)<\diam V_0$ для всех $v\in\Lambda\setminus\{0\}$
с $\|v\|^2<s^*$. Если $\Gamma\subset\Lambda$ задаёт правильную раскраску, то
\[
[\Lambda:\Gamma]\ \ge\ \frac{(s^*)^{n/2}}{\gamma_n^{n/2}\det\Lambda},
\]
где $\gamma_n$~--- постоянная Эрмита.
\end{proposition}

\begin{proof}
Правильность означает $D(v)\ge\diam V_0$ для всех $v\in\Gamma\setminus\{0\}$,
поэтому $\lambda_1(\Gamma)^2\ge s^*$; по неравенству Эрмита
$\lambda_1(\Gamma)^n\le\gamma_n^{n/2}\det\Gamma=\gamma_n^{n/2}[\Lambda:\Gamma]\det\Lambda$.
\end{proof}

По лемме~\ref{lem:inr} всегда годится $s^*=(\diam V_0+\lambda_1)^2$. Но нормы
векторов решётки пробегают не континуум, а \emph{оболочки}; если первая
оболочка выше этого порога запрещена целиком, $s^*$ поднимается скачком до
следующей. Запрет оболочки удостоверяется одной строкой.

\begin{lemma}[радиальный свидетель]\label{lem:rad}
Пусть $\langle\Lambda,\Lambda\rangle\subseteq\frac1e\Z$, $\|v\|^2=N$, и
рациональное $s\in(0,\frac12)$ таково, что для каждой оболочки
$\|u\|^2=M<4s^2N$ выполнено
$s\lfloor e\sqrt{NM}\rfloor/e\le M/2$. Тогда $sv\in V_0$ и
$D(v)\le(1-2s)\sqrt N$.
\end{lemma}

\begin{proof}
Надо проверить $s\langle v,u\rangle\le\|u\|^2/2$ для всех $u\in\Lambda$. При
$M=\|u\|^2\ge4s^2N$ это следует из неравенства Коши--Буняковского:
$s\langle v,u\rangle\le s\sqrt{NM}\le M/2$. При $M<4s^2N$ число
$\langle v,u\rangle$ кратно $1/e$ и не превосходит $\sqrt{NM}$, то есть не
превосходит $\lfloor e\sqrt{NM}\rfloor/e$, и работает условие леммы. Наконец,
$D(v)=2\dist(v/2,V_0)\le2\|v/2-sv\|=(1-2s)\sqrt N$.
\end{proof}

\begin{proof}[Доказательство теоремы~\ref{thm:e8}]
\emph{Решётка $E_8$:} $\lambda_1^2=2$, $R=1$, $\diam V_0=2$, $\det\Lambda=1$,
$\langle\Lambda,\Lambda\rangle=\Z$, нормы чётны. Правильность требует по
лемме~\ref{lem:inr} $\|v\|\ge2+\sqrt2$, то есть $\|v\|^2\ge6+4\sqrt2>11$,
и ближайшая оболочка~--- $12$. Пусть $\|v\|^2=12$; применим
лемму~\ref{lem:rad} с $e=1$, $s=1/4$: тогда $4s^2N=3$, отдельной проверки
требует лишь оболочка $M=2$, где $\lfloor\sqrt{24}\rfloor=4$ и
$\frac14\cdot4=1=M/2$. Значит, $D(v)\le\frac12\sqrt{12}=\sqrt3<2=\diam V_0$:
оболочка $12$ запрещена целиком, $s^*=14$, и предложение~\ref{prop:shell} с
$\gamma_8=2$ \cite{Blichfeldt} даёт $[\Lambda:\Gamma]\ge14^4/2^4=7^4=2401$.
Для $\Gamma=(3+\omega)E_8$ индекс равен $2401$, а раскраска правильна по
следствию~\ref{cor:widths}.

\emph{Решётка $A_3^*$} (ОЦК: $\Z^3\cup(\Z^3+(\frac12,\frac12,\frac12))$):
$\lambda_1^2=3/4$, $R^2=5/16$, $\det\Lambda=1/2$,
$\langle\Lambda,\Lambda\rangle\subseteq\frac14\Z$. Лемма~\ref{lem:inr} даёт
$\|v\|^2\ge(2R+\lambda_1)^2=2+\sqrt{15}/2>3{,}93$; нормы ОЦК~--- целые числа
и числа вида (сумма трёх нечётных квадратов)$/4$, так что в промежутке
$[3{,}93;\,5]$ лежат оболочки $4$, $19/4$, $5$. Для $\|v\|^2=4$ лемма~\ref{lem:rad} с
$e=4$, $s=1/4$: $4s^2N=1$, отдельно проверяется $M=3/4$, где
$\lfloor4\sqrt3\rfloor/4=3/2$ и $\frac14\cdot\frac32=\frac38=M/2$; значит,
$D(v)\le\frac12\cdot2=1<\sqrt5/2=\diam V_0$, и $s^*=19/4$. Предложение~\ref{prop:shell}
с $\gamma_3^{3/2}=\sqrt2$ \cite{ConwaySloane} даёт
$[\Lambda:\Gamma]^2\ge(19/4)^3\cdot2=6859/32>214$, то есть
$[\Lambda:\Gamma]\ge15$.
\end{proof}

Для ОЦК граница $15$ известна из общей теоремы Кулсона~\cite{Coulson}; новое
здесь~--- вывод из одного лишь спектра норм, который для $E_8$ даёт $2401$
вместо $511$. Тем же приёмом получаются полы $305$ на $E_6^*$ (рекорд $343$)
и $5\,688\,009\,064$ на решётке Лича (рекорд $7^{12}$; используется
$\gamma_{24}=4$~\cite{CohnKumar}). Отметим, что на решётках общего положения,
где нормы почти непрерывны, оболочечный пол не даёт выигрыша: минимальность
индекса на них остаётся открытым вопросом.

Программы, воспроизводящие численный контроль всех утверждений (в том
числе перебор всех векторов подрешёток в доказуемо полных окнах), открыты.\footnote{\texttt{https://github.com/ivaleo/Chromatic}}

\section*{ИСТОЧНИК ФИНАНСИРОВАНИЯ}

Работа выполнена без финансовой поддержки.

\section*{КОНФЛИКТ ИНТЕРЕСОВ}

Авторы заявляют об отсутствии конфликта интересов.

\begin{thebibliography}{99}

\bibitem{Soifer}
\textit{Soifer A.}
The Mathematical Coloring Book. New York: Springer, 2009. 607~p.

\bibitem{deGrey}
\textit{de Grey A.\,D.\,N.\,J.}
The chromatic number of the plane is at least 5 //
Geombinatorics. 2018. V.~28. №~1. P.~18--31.

\bibitem{Ivanov2011}
\textit{Иванов Л.\,Л.}
О хроматических числах $\R^2$ и $\R^3$ с интервалами запрещённых
расстояний // Вестник РУДН. Сер. Математика. Информатика. Физика. 2011.
№~1. С.~14--29.

\bibitem{ABPR}
\textit{Arman A., Bondarenko A., Prymak A., Radchenko D.}
Upper bounds on chromatic number of $\mathbb E^n$ in low dimensions //
Electron. J. Combin. 2024. V.~31. №~2. P2.35.
https://doi.org/10.37236/11440

\bibitem{ConwaySloane}
\textit{Conway J.\,H., Sloane N.\,J.\,A.}
Sphere Packings, Lattices and Groups. 3rd ed. New York: Springer, 1999.
703~p.

\bibitem{Coulson}
\textit{Coulson D.}
A 15-colouring of 3-space omitting distance one //
Discrete Math. 2002. V.~256. P.~83--90.

\bibitem{Blichfeldt}
\textit{Blichfeldt H.\,F.}
The minimum values of positive quadratic forms in six, seven and eight
variables // Math. Z. 1935. V.~39. P.~1--15.

\bibitem{CohnKumar}
\textit{Cohn H., Kumar A.}
Optimality and uniqueness of the Leech lattice among lattices //
Ann. of Math. (2). 2009. V.~170. №~3. P.~1003--1050.

\end{thebibliography}

\renewcommand\refname{References}

\begin{thebibliography}{99}

\bibitem{Soifer_e}
\textit{Soifer A.}
The Mathematical Coloring Book. New York: Springer, 2009. 607~p.

\bibitem{deGrey_e}
\textit{de Grey A.\,D.\,N.\,J.}
The chromatic number of the plane is at least 5 //
Geombinatorics. 2018. V.~28. No.~1. P.~18--31.

\bibitem{Ivanov2011_e}
\textit{Ivanov L.\,L.}
On the chromatic numbers of $\R^2$ and $\R^3$ with intervals of forbidden
distances // Bulletin of PFUR. Ser. Mathematics. Information Sciences.
Physics. 2011. No.~1. P.~14--29. (In Russian.)

\bibitem{ABPR_e}
\textit{Arman A., Bondarenko A., Prymak A., Radchenko D.}
Upper bounds on chromatic number of $\mathbb E^n$ in low dimensions //
Electron. J. Combin. 2024. V.~31. No.~2. P2.35.
https://doi.org/10.37236/11440

\bibitem{ConwaySloane_e}
\textit{Conway J.\,H., Sloane N.\,J.\,A.}
Sphere Packings, Lattices and Groups. 3rd ed. New York: Springer, 1999.
703~p.

\bibitem{Coulson_e}
\textit{Coulson D.}
A 15-colouring of 3-space omitting distance one //
Discrete Math. 2002. V.~256. P.~83--90.

\bibitem{Blichfeldt_e}
\textit{Blichfeldt H.\,F.}
The minimum values of positive quadratic forms in six, seven and eight
variables // Math. Z. 1935. V.~39. P.~1--15.

\bibitem{CohnKumar_e}
\textit{Cohn H., Kumar A.}
Optimality and uniqueness of the Leech lattice among lattices //
Ann. of Math. (2). 2009. V.~170. No.~3. P.~1003--1050.

\end{thebibliography}

\end{document}
